
\section{Erd\H{o}s Problem \#260}

\subsection*{1) ROUND-2 OBJECTIVE}

I pursue path \textbf{(C)} in the fail-safe sense: not a full proof of the original problem, but a strictly stronger theorem than Round~1 together with a sharper identification of the remaining obstruction.  Round~1 proved irrationality under the condition
\[
a_{n+1}-a_n-\lfloor \log_2 a_{n+1}\rfloor \longrightarrow \infty.
\]
The new goal is to weaken this to a mere \emph{limsup} condition.

\subsection*{2) Round-1 FOUNDATION USED}

I use only the Round~1 binary-support viewpoint:
\[
\frac{a_n}{2^{a_n}}
\]
contributes binary digits only in the interval
\[
[a_n-\lfloor \log_2 a_n\rfloor,\; a_n].
\]
Round~1 used eventual disjointness of these intervals to get a no-carry argument.

\subsection*{3) NEW INSIGHT / TOOL (ROUND-2)}

The new tool is a \emph{tail-stability lemma for binary digits}: if a dyadic partial sum has denominator $2^m$ and the remaining tail is less than $2^{-M}$, then the first $M-1$ binary digits are frozen, regardless of later carries.

This bypasses the Round~1 no-carry requirement entirely.  I only need infinitely many large gaps, not eventual gap separation.

\subsection*{4) ATTACK PLAN (ROUND-2)}

The exact Round~1 gap was this: eventual disjointness of digit blocks was used to control carries.  I replace that by a quantitative tail estimate.

Define
\[
L_n:=\lfloor \log_2 a_n\rfloor,\qquad
\Delta_n:=a_{n+1}-L_{n+1}-a_n-1.
\]
Thus $\Delta_n$ is the number of binary positions strictly between the last possible nonzero digit coming from the $n$th summand and the first possible nonzero digit coming from the $(n+1)$st summand.

The claim to prove is:

\medskip
\noindent
\textbf{New theorem.}
If
\[
\limsup_{n\to\infty}\Delta_n=\infty,
\]
then
\[
\sum_{n=1}^{\infty}\frac{a_n}{2^{a_n}}
\]
is irrational.
\medskip

This strictly strengthens the Round~1 theorem, which required $\Delta_n\to\infty$.

\subsection*{5) WORK (ROUND-2)}

Let
\[
S:=\sum_{n=1}^{\infty}\frac{a_n}{2^{a_n}},\qquad
P_n:=\sum_{k=1}^{n}\frac{a_k}{2^{a_k}},\qquad
T_n:=S-P_n=\sum_{k>n}\frac{a_k}{2^{a_k}}.
\]

\medskip
\noindent
\textbf{Step 1: a sharp upper bound for the tail.}

Since $a_{n+1}<a_{n+2}<\cdots$, we have
\[
T_n\leq \sum_{m=a_{n+1}}^{\infty}\frac{m}{2^m}.
\]
Using the standard identity
\[
\sum_{m=N}^{\infty}\frac{m}{2^m}=\frac{N+1}{2^{N-1}}
\qquad (N\geq 1),
\]
or equivalently
\[
\sum_{m=N+1}^{\infty}\frac{m}{2^m}=\frac{N+2}{2^N},
\]
we get
\[
T_n\leq \frac{a_{n+1}+2}{2^{a_{n+1}}}.
\]
Now $a_{n+1}<2^{L_{n+1}+1}$, hence $a_{n+1}+2<2^{L_{n+1}+2}$.  Therefore
\[
T_n<2^{L_{n+1}+2-a_{n+1}}.
\]
Because
\[
a_{n+1}=a_n+1+\Delta_n+L_{n+1},
\]
this becomes
\[
T_n<2^{-(a_n+\Delta_n-1)}.
\tag{260.1}
\]

\medskip
\noindent
\textbf{Step 2: long zero runs in the binary expansion of $S$.}

The partial sum $P_n$ is a dyadic rational with denominator dividing $2^{a_n}$.  Hence in its terminating binary expansion, all digits in positions strictly after $a_n$ are $0$.

From (260.1), adding $T_n$ can change only binary digits in positions at least $a_n+\Delta_n-1$.  Therefore the binary digits of $S$ in the positions
\[
a_n+1,\; a_n+2,\; \dots,\; a_n+\Delta_n-2
\]
agree with those of $P_n$, hence are all $0$.

So whenever $\Delta_n\geq 2$, the binary expansion of $S$ contains a zero run of length at least
\[
\Delta_n-2.
\]
Since $\limsup \Delta_n=\infty$, the binary expansion of $S$ contains arbitrarily long zero runs.

\medskip
\noindent
\textbf{Step 3: a rational number with arbitrarily long zero runs must be dyadic.}

Take the standard binary expansion, i.e. the one not ending in an infinite tail of $1$'s.  If a rational number is not dyadic, then its standard binary expansion is eventually periodic but not eventually zero.  If the eventual period has length $q$, then every sufficiently late block of $q$ consecutive digits contains at least one $1$; otherwise one whole period would be all zero, forcing termination.  Hence the lengths of zero runs are bounded.

Therefore a rational number whose standard binary expansion has arbitrarily long zero runs must be dyadic.

So, to finish, it remains only to rule out the possibility that $S$ is dyadic.

\medskip
\noindent
\textbf{Step 4: $S$ is not dyadic.}

Assume for contradiction that
\[
S=\frac{u}{2^M}
\]
for some integers $u,M\geq 0$.

Choose $n$ so large that $a_n\geq M$ and $\Delta_n\geq 2$; this is possible because $a_n\to\infty$ and $\limsup \Delta_n=\infty$.

Then $P_n$ has denominator dividing $2^{a_n}$, and so does $S$.  Hence
\[
T_n=S-P_n
\]
is a nonnegative multiple of $2^{-a_n}$.

But $T_n>0$ because there are infinitely many positive summands after the $n$th one, while by (260.1) and $\Delta_n\geq 2$ we have
\[
0<T_n<2^{-a_n}.
\]
This is impossible, because the only multiple of $2^{-a_n}$ lying in the interval $(0,2^{-a_n})$ is none at all.

This contradiction shows that $S$ is not dyadic.  Combined with Step~3, it follows that $S$ is irrational.

\medskip
\noindent
\textbf{Conclusion.}
If
\[
\limsup_{n\to\infty}\bigl(a_{n+1}-a_n-\lfloor \log_2 a_{n+1}\rfloor-1\bigr)=\infty,
\]
then
\[
\sum_{n=1}^{\infty}\frac{a_n}{2^{a_n}}
\]
is irrational.

\subsection*{6) ADVERSARIAL VERIFICATION}

\begin{itemize}
\item \textbf{Carry issue.}  The proof never assumes no-carry.  The only input is that the tail is numerically too small to alter early binary digits.
\item \textbf{Terminating-binary loophole.}  Arbitrarily long zero runs alone do not exclude dyadic rationals; this is why Step~4 is separate.
\item \textbf{Quantifiers.}  I only need infinitely many indices with large $\Delta_n$, not eventual control.  This is exactly why the hypothesis is a $\limsup$ condition.
\item \textbf{Strict advance over Round~1.}  Round~1 required $\Delta_n\to\infty$; the new theorem needs only $\limsup \Delta_n=\infty$.
\end{itemize}

\subsection*{7) FINAL (EXACTLY ONE)}

\textbf{UNRESOLVED (BUT STRICTLY ADVANCED)}

Round~2 proves a strictly stronger irrationality criterion than Round~1:
\[
\limsup_{n\to\infty}\bigl(a_{n+1}-a_n-\lfloor \log_2 a_{n+1}\rfloor-1\bigr)=\infty
\]
already suffices.  The genuine remaining obstruction is therefore concentrated in the regime where
\[
a_{n+1}-a_n-\lfloor \log_2 a_{n+1}\rfloor
\]
stays bounded above, so the tail can repeatedly interact with the earlier binary digits.

\subsection*{8) COMPLETION ESTIMATE}

COMPLETION: 50\%

\subsection*{9) REFERENCES}

\begin{itemize}
\item [EP260] T. F. Bloom, \emph{Erd\H{o}s Problem \#260}, problem page, accessed 2026-03-07.
\item [Er88c] P. Erd\H{o}s, \emph{On the irrationality of certain series: problems and results}, 1988; see p.~103 for Problem \#260 and the discussion of stronger sufficient hypotheses.
\end{itemize}

\bigskip
\hrule
\bigskip

